% 子文件：6.无穷级数.tex
% 对应教材：同济《高等数学》下册 第十二章 无穷级数
% 节次：§12.1 常数项级数的概念和性质｜§12.2 常数项级数的审敛法｜§12.3 幂级数
%       §12.4 函数展开成幂级数｜§12.5 函数的幂级数展开式的应用
%       §12.6 函数项级数的一致收敛性｜§12.7 傅里叶级数｜§12.8 一般周期函数的傅里叶级数
% 主文档使用 % 子文件：6.无穷级数.tex
% 对应教材：同济《高等数学》下册 第十二章 无穷级数
% 节次：§12.1 常数项级数的概念和性质｜§12.2 常数项级数的审敛法｜§12.3 幂级数
%       §12.4 函数展开成幂级数｜§12.5 函数的幂级数展开式的应用
%       §12.6 函数项级数的一致收敛性｜§12.7 傅里叶级数｜§12.8 一般周期函数的傅里叶级数
% 主文档使用 % 子文件：6.无穷级数.tex
% 对应教材：同济《高等数学》下册 第十二章 无穷级数
% 节次：§12.1 常数项级数的概念和性质｜§12.2 常数项级数的审敛法｜§12.3 幂级数
%       §12.4 函数展开成幂级数｜§12.5 函数的幂级数展开式的应用
%       §12.6 函数项级数的一致收敛性｜§12.7 傅里叶级数｜§12.8 一般周期函数的傅里叶级数
% 主文档使用 % 子文件：6.无穷级数.tex
% 对应教材：同济《高等数学》下册 第十二章 无穷级数
% 节次：§12.1 常数项级数的概念和性质｜§12.2 常数项级数的审敛法｜§12.3 幂级数
%       §12.4 函数展开成幂级数｜§12.5 函数的幂级数展开式的应用
%       §12.6 函数项级数的一致收敛性｜§12.7 傅里叶级数｜§12.8 一般周期函数的傅里叶级数
% 主文档使用 \input{6.无穷级数} 引入，不含 document 环境

\section{常数项级数的概念和性质}
\subsection{概念}

\textbf{定义}\quad 给定数列 $u_1,u_2,\cdots,u_n,\cdots$，称 $\sum\limits_{n=1}^{\infty}u_n=u_1+u_2+\cdots+u_n+\cdots$ 为（常数项）\textbf{无穷级数}，$u_n$ 为一般项。作部分和 $s_n=u_1+\cdots+u_n$，若 $\lim\limits_{n\to\infty}s_n=s$（常数），称级数\textbf{收敛}，$s$ 为\textbf{和}；若 $\{s_n\}$ 无极限，称级数\textbf{发散}。收敛时 $r_n=s-s_n=\sum\limits_{k=n+1}^{\infty}u_k$ 称\textbf{余项}，且 $r_n\to0$。

\textbf{两个基准级数}（必须记牢）：
\keyformula{\sum_{n=0}^{\infty}aq^{n}=\frac{a}{1-q}\ (|q|<1),\quad \text{发散}\ (|q|\ge1);\qquad \sum_{n=1}^{\infty}\frac{1}{n^{p}}\ \text{收敛}\iff p>1}
几何级数：$|q|<1$ 收敛于 $\frac{a}{1-q}$；$p$ 级数：$p>1$ 收敛，$p\le1$ 发散（$p=1$ 为调和级数，发散）。

\subsection{收敛级数的基本性质}

\begin{enumerate}
  \item \textbf{线性}：$\sum u_n=s,\ \sum v_n=t\Rightarrow\sum(\alpha u_n+\beta v_n)=\alpha s+\beta t$。
  \item \textbf{去掉、加上或改变有限项}不改变收敛性（收敛时和一般会变）。
  \item \textbf{收敛的必要条件}：若 $\sum u_n$ 收敛，则 $\lim\limits_{n\to\infty}u_n=0$。
  \item \textbf{收敛级数可任意加括号}（和不变）；但\textbf{去括号}不一定成立。
  \item \textbf{柯西审敛原理}：$\sum u_n$ 收敛 $\iff$ 对任意 $\varepsilon>0$ 存在 $N$，当 $n>N$ 时对任意正整数 $p$ 有 $|u_{n+1}+\cdots+u_{n+p}|<\varepsilon$。
\end{enumerate}

\corepoint{性质 3 是\textbf{必要}条件，只用来判\textbf{发散}：$\lim u_n\ne0$ 则必发散；但 $\lim u_n=0$ 推不出收敛（调和级数 $\sum\frac1n$ 即反例），这是判敛第一步。性质 4 的单向性常考：加括号收敛 $\not\Rightarrow$ 原级数收敛。}

\section{常数项级数的审敛法}
\subsection{正项级数及其审敛法}

若 $u_n\ge0$，称 $\sum u_n$ 为\textbf{正项级数}。

\textbf{定理 1（基本定理）}\quad 正项级数收敛的\textbf{充要条件}是部分和数列 $\{s_n\}$ \textbf{有界}（$s_n$ 单调增，单调有界必收敛）。

\textbf{定理 2（比较审敛法）}\quad 设 $\sum u_n,\sum v_n$ 为正项级数且 $u_n\le v_n$：若 $\sum v_n$ 收敛则 $\sum u_n$ 收敛；若 $\sum u_n$ 发散则 $\sum v_n$ 发散。

\textbf{推论（极限形式）}\quad 设 $\lim\limits_{n\to\infty}\frac{u_n}{v_n}=l$：$0<l<+\infty$ 时两者\textbf{同敛散}；$l=0$ 且 $\sum v_n$ 收敛则 $\sum u_n$ 收敛；$l=+\infty$ 且 $\sum v_n$ 发散则 $\sum u_n$ 发散。

\textbf{定理 3（比值审敛法，达朗贝尔）}\quad 对正项级数 $\sum u_n$，若 $\lim\limits_{n\to\infty}\frac{u_{n+1}}{u_n}=\rho$，则 $\rho<1$ 收敛，$\rho>1$（含 $+\infty$）发散，$\rho=1$ \textbf{失效}。

\textbf{定理 4（根值审敛法，柯西）}\quad 若 $\lim\limits_{n\to\infty}\sqrt[n]{u_n}=\rho$，结论同定理 3。

\textbf{定理 5（积分审敛法）}\quad 设 $f(x)$ 在 $[1,+\infty)$ 上连续、非负、单调减少且 $f(n)=u_n$，则 $\sum\limits_{n=1}^{\infty}u_n$ 与 $\int_1^{+\infty}f(x)\,\mathrm{d}x$ \textbf{同时收敛或同时发散}。

\corepoint{比较法的\textbf{关键是选比较对象}：常用几何级数、$p$ 级数、$\frac{1}{n\ln^pn}$ 型。选法看通项的阶：含 $n!$、$a^n$、$n^n$ 优先用\textbf{比值}；通项为 $(f(n))^n$ 形优先用\textbf{根值}；$n$ 的有理分式（含根号）优先用\textbf{比较法}（与 $\frac{1}{n^p}$ 比，看分母最高次与分子之差）。比值、根值法在 $\rho=1$ 时\textbf{不能用}（$\sum\frac1n$、$\sum\frac{1}{n^2}$ 都给出 $\rho=1$），这是最常见失分点。}

\subsection{交错级数及其审敛法}

形如 $\sum\limits_{n=1}^{\infty}(-1)^{n-1}u_n=u_1-u_2+u_3-\cdots$（$u_n>0$）的级数称\textbf{交错级数}。

\textbf{定理 6（莱布尼茨判别法）}\quad 若 (1) $u_n\ge u_{n+1}$（$n=1,2,\cdots$）；(2) $\lim\limits_{n\to\infty}u_n=0$，则 $\sum(-1)^{n-1}u_n$ \textbf{收敛}，且 $s\le u_1$，余项 $|r_n|\le u_{n+1}$。

\corepoint{两条件缺一不可：$u_n$ \textbf{单调不增}且 $\to0$；只看 $\to0$ 不够（不单调时可能发散）。非交错级数不能套用，应先拆项。余项估计 $|r_n|\le u_{n+1}$ 常用于近似计算定误差。}

\subsection{绝对收敛与条件收敛}

对任意项级数 $\sum u_n$：若 $\sum|u_n|$ 收敛，称 $\sum u_n$ \textbf{绝对收敛}；若 $\sum u_n$ 收敛而 $\sum|u_n|$ 发散，称\textbf{条件收敛}。

\textbf{定理 7}\quad 绝对收敛 $\Rightarrow$ 收敛。

\textbf{定理 8（绝对收敛级数的性质）}\quad 绝对收敛级数可\textbf{任意交换项次序}而不改变其和；两个绝对收敛级数的\textbf{柯西乘积} $\sum\limits_{n=1}^{\infty}\left(\sum\limits_{k=1}^{n}u_kv_{n-k+1}\right)$ 也绝对收敛，和为 $(\sum u_n)(\sum v_n)$。条件收敛级数不能随意重排。

\corepoint{判任意项级数的\textbf{标准流程}：先看 $\sum|u_n|$——收敛即\textbf{绝对收敛}，结束；若发散，再判原级数（交错用莱布尼茨，一般用定义）——收敛则为\textbf{条件收敛}。$|u_n|$ 发散\textbf{不能}推出 $u_n$ 发散，这是最易搞反的一步。}

\section{幂级数}
\subsection{函数项级数}

设 $u_n(x)$ 定义在区间 $I$ 上，称 $\sum\limits_{n=1}^{\infty}u_n(x)$ 为\textbf{函数项级数}。若 $\sum u_n(x_0)$ 收敛，称 $x_0$ 为收敛点，收敛点全体称\textbf{收敛域}，其上的 $s(x)=\sum u_n(x)$ 称\textbf{和函数}。

\subsection{幂级数及其收敛性}

形如 $\sum\limits_{n=0}^{\infty}a_n(x-x_0)^n$ 的级数称\textbf{幂级数}（$x_0=0$ 时即 $\sum a_nx^n$）。

\textbf{阿贝尔定理}\quad 若 $\sum a_nx^n$ 在 $x=x_0\ne0$ 处收敛，则对 $|x|<|x_0|$ 的一切 $x$ \textbf{绝对收敛}；若在 $x_0$ 处发散，则对 $|x|>|x_0|$ 的一切 $x$ \textbf{发散}。

故收敛性只有三种：仅在 $x=0$ 收敛（$R=0$）；全数轴收敛（$R=+\infty$）；存在 $R>0$，$|x|<R$ 绝对收敛、$|x|>R$ 发散——$R$ 称\textbf{收敛半径}，$(-R,R)$ 称\textbf{收敛区间}，端点 $x=\pm R$ 须单独判定，最终得\textbf{收敛域}。

\textbf{收敛半径的求法}：对 $\sum a_nx^n$，若
\keyformula{\rho=\lim_{n\to\infty}\left|\frac{a_{n+1}}{a_n}\right|\quad\text{或}\quad \rho=\lim_{n\to\infty}\sqrt[n]{|a_n|},\qquad R=\frac{1}{\rho}}
（约定 $\rho=0$ 时 $R=+\infty$；$\rho=+\infty$ 时 $R=0$。）

\corepoint{求收敛域\textbf{三步}：(1) 用公式求 $R$（$a_n$ 是 $x^n$ 的系数；缺项级数如 $\sum a_nx^{2n}$ 不能直接套公式，应把 $x^2$ 当整体用比值法）；(2) 得 $(-R,R)$；(3) \textbf{端点单独讨论}——代入 $x=\pm R$ 化成数项级数再判，两端可能一敛一散。只写开区间不讨论端点是最常见的扣分点。}

\subsection{幂级数的运算}

设 $\sum a_nx^n=f(x)$，$\sum b_nx^n=g(x)$，收敛半径 $R_1,R_2$，$R=\min\{R_1,R_2\}$，则在 $(-R,R)$ 内
\[
\sum a_nx^n\pm\sum b_nx^n=\sum(a_n\pm b_n)x^n,\qquad
\left(\sum a_nx^n\right)\left(\sum b_nx^n\right)=\sum_{n=0}^{\infty}\left(\sum_{k=0}^{n}a_kb_{n-k}\right)x^n .
\]
在 $(-R,R)$ 内和函数连续，且可\textbf{逐项求导}、\textbf{逐项积分}：
\keyformula{s'(x)=\sum_{n=1}^{\infty}na_nx^{n-1},\qquad \int_0^{x}s(t)\,\mathrm{d}t=\sum_{n=0}^{\infty}\frac{a_n}{n+1}x^{n+1}}
运算后\textbf{收敛半径不变}，但\textbf{端点处的收敛性可能改变}。

\corepoint{逐项求导降次、可能"丢失"常数项，逐项积分升次并添 $x$ 因子——求\textbf{和函数}正靠这两招把系数凑成可积分/可求导的形状，再回到已知展开式。逐项运算只在\textbf{开区间}内保证成立，端点要重新讨论。}

\section{函数展开成幂级数}
\subsection{泰勒级数与展开条件}

若 $f(x)$ 在 $x_0$ 的某邻域内有任意阶导数，则称
\[
\sum_{n=0}^{\infty}\frac{f^{(n)}(x_0)}{n!}(x-x_0)^n=f(x_0)+f'(x_0)(x-x_0)+\frac{f''(x_0)}{2!}(x-x_0)^2+\cdots
\]
为 $f(x)$ 在 $x_0$ 处的\textbf{泰勒级数}；$x_0=0$ 时称\textbf{麦克劳林级数}。

\textbf{定理（展开的充要条件）}\quad $f(x)$ 在该邻域内能展开成泰勒级数 $\iff$ 泰勒公式的余项满足 $\lim\limits_{n\to\infty}R_n(x)=0$。

\textbf{定理（唯一性）}\quad 若 $f(x)$ 能展开成 $x-x_0$ 的幂级数，则系数必为 $a_n=\dfrac{f^{(n)}(x_0)}{n!}$，展开式唯一。

\subsection{五个基本展开式}

\begin{center}
\begin{tabular}{@{}lll@{}}
\toprule
函数 & 展开式 & 收敛域 \\
\midrule
$\mathrm{e}^{x}$ & $\sum\limits_{n=0}^{\infty}\dfrac{x^{n}}{n!}=1+x+\dfrac{x^{2}}{2!}+\cdots$ & $(-\infty,+\infty)$ \\
$\sin x$ & $\sum\limits_{n=0}^{\infty}\dfrac{(-1)^{n}x^{2n+1}}{(2n+1)!}=x-\dfrac{x^{3}}{3!}+\dfrac{x^{5}}{5!}-\cdots$ & $(-\infty,+\infty)$ \\
$\cos x$ & $\sum\limits_{n=0}^{\infty}\dfrac{(-1)^{n}x^{2n}}{(2n)!}=1-\dfrac{x^{2}}{2!}+\dfrac{x^{4}}{4!}-\cdots$ & $(-\infty,+\infty)$ \\
$\ln(1+x)$ & $\sum\limits_{n=1}^{\infty}\dfrac{(-1)^{n-1}x^{n}}{n}=x-\dfrac{x^{2}}{2}+\dfrac{x^{3}}{3}-\cdots$ & $(-1,1]$ \\
$(1+x)^{m}$ & $1+mx+\dfrac{m(m-1)}{2!}x^{2}+\cdots$（$m$ 非非负整数） & $(-1,1)$ \\
\bottomrule
\end{tabular}
\end{center}

\keyformula{\mathrm{e}^{x}=\sum_{n=0}^{\infty}\frac{x^{n}}{n!},\qquad \frac{1}{1-x}=\sum_{n=0}^{\infty}x^{n}\ (|x|<1)}

\subsection{间接展开法}

直接按 $a_n=\frac{f^{(n)}(0)}{n!}$ 求系数（直接展开法）计算量大，通常用\textbf{间接展开法}：由已知展开式经\textbf{变量代换}、四则运算、\textbf{逐项求导}、\textbf{逐项积分}得新展开式，由唯一性保证正确。

\textbf{常用套路}：(1) 把 $f$ 化成 $\mathrm{e}^{u},\sin u,\ln(1+u),(1+u)^m,\frac{1}{1-u}$ 的形式，令 $u=\varphi(x)$；(2) 有理分式先\textbf{拆成部分分式}再用 $\frac{1}{1\mp x}$；(3) 含 $\arctan x,\arcsin x,\ln(1+x)$ 的先求导化成有理式，展开后再逐项积分。

\corepoint{间接展开最易错的是\textbf{收敛域}：代换后范围要跟着换。由 $\frac{1}{1-x}=\sum x^n\,(|x|<1)$ 得 $\frac{1}{1+x^2}=\sum(-1)^nx^{2n}$，收敛域仍是 $|x|<1$；而 $\ln(1+x)$ 逐项积分后在 $x=1$ 也收敛（得 $(-1,1]$）。\textbf{端点是否包含必须重新判}，不能照抄原级数区间。}

\section{函数的幂级数展开式的应用}

利用幂级数可作近似计算，误差由\textbf{余项}估计：(1) \textbf{求函数近似值}——把 $x$ 取小，展开到所需精度，由余项定项数；(2) \textbf{求积分近似值}——被积函数展开后\textbf{逐项积分}再截断；(3) \textbf{求极限}——用展开式代替函数，同阶项相消。

\textbf{微分方程的幂级数解法}：对 $y'=f(x,y)$，若 $f$ 在 $(x_0,y_0)$ 附近可展成幂级数，可设 $y=\sum a_n(x-x_0)^n$，代入比较同次幂系数逐次定出 $a_n$。

\corepoint{核心是\textbf{误差控制}：问"精确到 $10^{-4}$"就要由 $|r_n|\le10^{-4}$ 反解 $n$。交错级数（满足莱布尼茨条件）的余项不超过第一个被舍去项，是最快的估计。}

\section{函数项级数的一致收敛性}

\textbf{定义}\quad 设 $\sum u_n(x)$ 在 $I$ 上收敛于 $s(x)$。若对任意 $\varepsilon>0$，存在\textbf{与 $x$ 无关}的 $N$，使 $n>N$ 时对一切 $x\in I$ 有 $|r_n(x)|<\varepsilon$，则称它在 $I$ 上\textbf{一致收敛}。\textbf{魏尔斯特拉斯判别法}：若 $|u_n(x)|\le M_n$ 且 $\sum M_n$ 收敛，则 $\sum u_n(x)$ 一致收敛。一致收敛级数具有\textbf{连续性}、可\textbf{逐项积分}、可\textbf{逐项求导}（导级数须一致收敛）。

\corepoint{与普通收敛的差别只在 $N$ 是否\textbf{与 $x$ 无关}；有了它才能逐项求极限、积分、求导。幂级数在其收敛区间内部一致收敛，这正是它可逐项运算的依据。}

\section{傅里叶级数}
\subsection{三角级数与三角函数系的正交性}

形如 $\frac{a_0}{2}+\sum\limits_{n=1}^{\infty}(a_n\cos nx+b_n\sin nx)$ 的级数称\textbf{三角级数}。三角函数系 $1,\cos x,\sin x,\cos 2x,\sin 2x,\cdots$ 在 $[-\pi,\pi]$ 上\textbf{正交}：
\[
\int_{-\pi}^{\pi}\cos mx\cos nx\,\mathrm{d}x=\int_{-\pi}^{\pi}\sin mx\sin nx\,\mathrm{d}x=\int_{-\pi}^{\pi}\cos mx\sin nx\,\mathrm{d}x=0\quad(m\ne n),
\]
而 $\int_{-\pi}^{\pi}1^2\mathrm{d}x=2\pi$，$\int_{-\pi}^{\pi}\cos^2nx\,\mathrm{d}x=\int_{-\pi}^{\pi}\sin^2nx\,\mathrm{d}x=\pi$。

\subsection{函数展开成傅里叶级数}

设 $f(x)$ 是周期为 $2\pi$ 的周期函数，由正交性逐项积分得\textbf{欧拉—傅里叶公式}
\keyformula{a_n=\frac{1}{\pi}\int_{-\pi}^{\pi}f(x)\cos nx\,\mathrm{d}x\ (n=0,1,2,\cdots),\qquad b_n=\frac{1}{\pi}\int_{-\pi}^{\pi}f(x)\sin nx\,\mathrm{d}x\ (n=1,2,\cdots)}
由此系数确定的三角级数称 $f$ 的\textbf{傅里叶级数}。

\textbf{狄利克雷收敛定理}\quad 设 $f$ 是周期 $2\pi$ 的周期函数，若 (1) 一个周期内\textbf{连续或只有有限个第一类间断点}；(2) 一个周期内\textbf{至多有有限个极值点}，则其傅里叶级数收敛，且
\[
\frac{a_0}{2}+\sum_{n=1}^{\infty}(a_n\cos nx+b_n\sin nx)=
\begin{cases}
f(x), & x\ \text{为连续点},\\[3pt]
\dfrac{f(x^-)+f(x^+)}{2}, & x\ \text{为间断点}.
\end{cases}
\]

\corepoint{两个要点：(1) 系数公式的 $\frac{1}{\pi}$ 与 $a_0$ 前的 $\frac12$ 配套，写错全错；(2) 收敛定理给的是"和函数"——\textbf{间断点}处收敛于左右极限的\textbf{平均值}，端点 $x=\pm\pi$ 处收敛于 $\frac{f(-\pi^+)+f(\pi^-)}{2}$。问"展开式在某点等于什么"，往往就考这个平均值。}

\subsection{正弦级数和余弦级数}

\begin{itemize}
  \item $f$ 为\textbf{奇函数}：$a_n=0$，$b_n=\frac{2}{\pi}\int_0^{\pi}f(x)\sin nx\,\mathrm{d}x$，级数只含正弦项（\textbf{正弦级数}）。
  \item $f$ 为\textbf{偶函数}：$b_n=0$，$a_n=\frac{2}{\pi}\int_0^{\pi}f(x)\cos nx\,\mathrm{d}x$，级数只含余弦项（\textbf{余弦级数}）。
\end{itemize}

\textbf{有限区间上的展开（奇延拓与偶延拓）}：若 $f$ 只在 $[0,\pi]$ 上有定义，要展成正弦级数就作\textbf{奇延拓}（$F(-x)=-F(x)$）再按奇函数求 $b_n$；要展成余弦级数就作\textbf{偶延拓}（$F(-x)=F(x)$）再按偶函数求 $a_n$。所得级数在 $(0,\pi)$ 内收敛于 $f(x)$。

\corepoint{奇延拓的正弦级数在 $x=0$ 处收敛于 $0$（除非 $f(0)=0$）；偶延拓的余弦级数在 $x=0$ 处收敛于 $f(0^+)$。这是"端点值"类题目的标准考法；延拓只是求系数的技巧，原函数在 $[0,\pi]$ 上并未改变。}

\section{一般周期函数的傅里叶级数}
\subsection{周期为 $2l$ 的周期函数}

设 $f(x)$ 的周期为 $2l$，令 $t=\frac{\pi x}{l}$ 可化为周期 $2\pi$ 的情形，得
\keyformula{f(x)\sim\frac{a_0}{2}+\sum_{n=1}^{\infty}\left(a_n\cos\frac{n\pi x}{l}+b_n\sin\frac{n\pi x}{l}\right)}
其中
\[
a_n=\frac{1}{l}\int_{-l}^{l}f(x)\cos\frac{n\pi x}{l}\,\mathrm{d}x\ (n=0,1,2,\cdots),\qquad
b_n=\frac{1}{l}\int_{-l}^{l}f(x)\sin\frac{n\pi x}{l}\,\mathrm{d}x\ (n=1,2,\cdots).
\]
$f$ 为奇函数时 $a_n=0$，$b_n=\frac{2}{l}\int_0^{l}f\sin\frac{n\pi x}{l}\mathrm{d}x$；为偶函数时 $b_n=0$，$a_n=\frac{2}{l}\int_0^{l}f\cos\frac{n\pi x}{l}\mathrm{d}x$。狄利克雷定理同样成立。

\subsection{傅里叶级数的复数形式（简）}

由欧拉公式 $\mathrm{e}^{\mathrm{i}\theta}=\cos\theta+\mathrm{i}\sin\theta$ 得复数形式
\keyformula{f(x)\sim\sum_{n=-\infty}^{+\infty}c_n\mathrm{e}^{\mathrm{i}\frac{n\pi x}{l}},\qquad c_n=\frac{1}{2l}\int_{-l}^{l}f(x)\,\mathrm{e}^{-\mathrm{i}\frac{n\pi x}{l}}\,\mathrm{d}x\ (n=0,\pm1,\pm2,\cdots)}
周期为 $2\pi$ 时 $c_n=\frac{1}{2\pi}\int_{-\pi}^{\pi}f(x)\mathrm{e}^{-\mathrm{i}nx}\mathrm{d}x$，且 $c_0=\frac{a_0}{2}$，$c_n=\frac{a_n-\mathrm{i}b_n}{2}$，$c_{-n}=\frac{a_n+\mathrm{i}b_n}{2}$。

\corepoint{周期 $2l$ 的记法：把 $2\pi$ 情形的 $\cos nx$ 换成 $\cos\frac{n\pi x}{l}$，系数前的 $\frac{1}{\pi}$ 换成 $\frac{1}{l}$（区间长度由 $2\pi$ 变 $2l$）。\textbf{积分区间须与周期一致}（$[-l,l]$ 或任一长 $2l$ 的区间），否则系数全错。}
 引入，不含 document 环境

\section{常数项级数的概念和性质}
\subsection{概念}

\textbf{定义}\quad 给定数列 $u_1,u_2,\cdots,u_n,\cdots$，称 $\sum\limits_{n=1}^{\infty}u_n=u_1+u_2+\cdots+u_n+\cdots$ 为（常数项）\textbf{无穷级数}，$u_n$ 为一般项。作部分和 $s_n=u_1+\cdots+u_n$，若 $\lim\limits_{n\to\infty}s_n=s$（常数），称级数\textbf{收敛}，$s$ 为\textbf{和}；若 $\{s_n\}$ 无极限，称级数\textbf{发散}。收敛时 $r_n=s-s_n=\sum\limits_{k=n+1}^{\infty}u_k$ 称\textbf{余项}，且 $r_n\to0$。

\textbf{两个基准级数}（必须记牢）：
\keyformula{\sum_{n=0}^{\infty}aq^{n}=\frac{a}{1-q}\ (|q|<1),\quad \text{发散}\ (|q|\ge1);\qquad \sum_{n=1}^{\infty}\frac{1}{n^{p}}\ \text{收敛}\iff p>1}
几何级数：$|q|<1$ 收敛于 $\frac{a}{1-q}$；$p$ 级数：$p>1$ 收敛，$p\le1$ 发散（$p=1$ 为调和级数，发散）。

\subsection{收敛级数的基本性质}

\begin{enumerate}
  \item \textbf{线性}：$\sum u_n=s,\ \sum v_n=t\Rightarrow\sum(\alpha u_n+\beta v_n)=\alpha s+\beta t$。
  \item \textbf{去掉、加上或改变有限项}不改变收敛性（收敛时和一般会变）。
  \item \textbf{收敛的必要条件}：若 $\sum u_n$ 收敛，则 $\lim\limits_{n\to\infty}u_n=0$。
  \item \textbf{收敛级数可任意加括号}（和不变）；但\textbf{去括号}不一定成立。
  \item \textbf{柯西审敛原理}：$\sum u_n$ 收敛 $\iff$ 对任意 $\varepsilon>0$ 存在 $N$，当 $n>N$ 时对任意正整数 $p$ 有 $|u_{n+1}+\cdots+u_{n+p}|<\varepsilon$。
\end{enumerate}

\corepoint{性质 3 是\textbf{必要}条件，只用来判\textbf{发散}：$\lim u_n\ne0$ 则必发散；但 $\lim u_n=0$ 推不出收敛（调和级数 $\sum\frac1n$ 即反例），这是判敛第一步。性质 4 的单向性常考：加括号收敛 $\not\Rightarrow$ 原级数收敛。}

\section{常数项级数的审敛法}
\subsection{正项级数及其审敛法}

若 $u_n\ge0$，称 $\sum u_n$ 为\textbf{正项级数}。

\textbf{定理 1（基本定理）}\quad 正项级数收敛的\textbf{充要条件}是部分和数列 $\{s_n\}$ \textbf{有界}（$s_n$ 单调增，单调有界必收敛）。

\textbf{定理 2（比较审敛法）}\quad 设 $\sum u_n,\sum v_n$ 为正项级数且 $u_n\le v_n$：若 $\sum v_n$ 收敛则 $\sum u_n$ 收敛；若 $\sum u_n$ 发散则 $\sum v_n$ 发散。

\textbf{推论（极限形式）}\quad 设 $\lim\limits_{n\to\infty}\frac{u_n}{v_n}=l$：$0<l<+\infty$ 时两者\textbf{同敛散}；$l=0$ 且 $\sum v_n$ 收敛则 $\sum u_n$ 收敛；$l=+\infty$ 且 $\sum v_n$ 发散则 $\sum u_n$ 发散。

\textbf{定理 3（比值审敛法，达朗贝尔）}\quad 对正项级数 $\sum u_n$，若 $\lim\limits_{n\to\infty}\frac{u_{n+1}}{u_n}=\rho$，则 $\rho<1$ 收敛，$\rho>1$（含 $+\infty$）发散，$\rho=1$ \textbf{失效}。

\textbf{定理 4（根值审敛法，柯西）}\quad 若 $\lim\limits_{n\to\infty}\sqrt[n]{u_n}=\rho$，结论同定理 3。

\textbf{定理 5（积分审敛法）}\quad 设 $f(x)$ 在 $[1,+\infty)$ 上连续、非负、单调减少且 $f(n)=u_n$，则 $\sum\limits_{n=1}^{\infty}u_n$ 与 $\int_1^{+\infty}f(x)\,\mathrm{d}x$ \textbf{同时收敛或同时发散}。

\corepoint{比较法的\textbf{关键是选比较对象}：常用几何级数、$p$ 级数、$\frac{1}{n\ln^pn}$ 型。选法看通项的阶：含 $n!$、$a^n$、$n^n$ 优先用\textbf{比值}；通项为 $(f(n))^n$ 形优先用\textbf{根值}；$n$ 的有理分式（含根号）优先用\textbf{比较法}（与 $\frac{1}{n^p}$ 比，看分母最高次与分子之差）。比值、根值法在 $\rho=1$ 时\textbf{不能用}（$\sum\frac1n$、$\sum\frac{1}{n^2}$ 都给出 $\rho=1$），这是最常见失分点。}

\subsection{交错级数及其审敛法}

形如 $\sum\limits_{n=1}^{\infty}(-1)^{n-1}u_n=u_1-u_2+u_3-\cdots$（$u_n>0$）的级数称\textbf{交错级数}。

\textbf{定理 6（莱布尼茨判别法）}\quad 若 (1) $u_n\ge u_{n+1}$（$n=1,2,\cdots$）；(2) $\lim\limits_{n\to\infty}u_n=0$，则 $\sum(-1)^{n-1}u_n$ \textbf{收敛}，且 $s\le u_1$，余项 $|r_n|\le u_{n+1}$。

\corepoint{两条件缺一不可：$u_n$ \textbf{单调不增}且 $\to0$；只看 $\to0$ 不够（不单调时可能发散）。非交错级数不能套用，应先拆项。余项估计 $|r_n|\le u_{n+1}$ 常用于近似计算定误差。}

\subsection{绝对收敛与条件收敛}

对任意项级数 $\sum u_n$：若 $\sum|u_n|$ 收敛，称 $\sum u_n$ \textbf{绝对收敛}；若 $\sum u_n$ 收敛而 $\sum|u_n|$ 发散，称\textbf{条件收敛}。

\textbf{定理 7}\quad 绝对收敛 $\Rightarrow$ 收敛。

\textbf{定理 8（绝对收敛级数的性质）}\quad 绝对收敛级数可\textbf{任意交换项次序}而不改变其和；两个绝对收敛级数的\textbf{柯西乘积} $\sum\limits_{n=1}^{\infty}\left(\sum\limits_{k=1}^{n}u_kv_{n-k+1}\right)$ 也绝对收敛，和为 $(\sum u_n)(\sum v_n)$。条件收敛级数不能随意重排。

\corepoint{判任意项级数的\textbf{标准流程}：先看 $\sum|u_n|$——收敛即\textbf{绝对收敛}，结束；若发散，再判原级数（交错用莱布尼茨，一般用定义）——收敛则为\textbf{条件收敛}。$|u_n|$ 发散\textbf{不能}推出 $u_n$ 发散，这是最易搞反的一步。}

\section{幂级数}
\subsection{函数项级数}

设 $u_n(x)$ 定义在区间 $I$ 上，称 $\sum\limits_{n=1}^{\infty}u_n(x)$ 为\textbf{函数项级数}。若 $\sum u_n(x_0)$ 收敛，称 $x_0$ 为收敛点，收敛点全体称\textbf{收敛域}，其上的 $s(x)=\sum u_n(x)$ 称\textbf{和函数}。

\subsection{幂级数及其收敛性}

形如 $\sum\limits_{n=0}^{\infty}a_n(x-x_0)^n$ 的级数称\textbf{幂级数}（$x_0=0$ 时即 $\sum a_nx^n$）。

\textbf{阿贝尔定理}\quad 若 $\sum a_nx^n$ 在 $x=x_0\ne0$ 处收敛，则对 $|x|<|x_0|$ 的一切 $x$ \textbf{绝对收敛}；若在 $x_0$ 处发散，则对 $|x|>|x_0|$ 的一切 $x$ \textbf{发散}。

故收敛性只有三种：仅在 $x=0$ 收敛（$R=0$）；全数轴收敛（$R=+\infty$）；存在 $R>0$，$|x|<R$ 绝对收敛、$|x|>R$ 发散——$R$ 称\textbf{收敛半径}，$(-R,R)$ 称\textbf{收敛区间}，端点 $x=\pm R$ 须单独判定，最终得\textbf{收敛域}。

\textbf{收敛半径的求法}：对 $\sum a_nx^n$，若
\keyformula{\rho=\lim_{n\to\infty}\left|\frac{a_{n+1}}{a_n}\right|\quad\text{或}\quad \rho=\lim_{n\to\infty}\sqrt[n]{|a_n|},\qquad R=\frac{1}{\rho}}
（约定 $\rho=0$ 时 $R=+\infty$；$\rho=+\infty$ 时 $R=0$。）

\corepoint{求收敛域\textbf{三步}：(1) 用公式求 $R$（$a_n$ 是 $x^n$ 的系数；缺项级数如 $\sum a_nx^{2n}$ 不能直接套公式，应把 $x^2$ 当整体用比值法）；(2) 得 $(-R,R)$；(3) \textbf{端点单独讨论}——代入 $x=\pm R$ 化成数项级数再判，两端可能一敛一散。只写开区间不讨论端点是最常见的扣分点。}

\subsection{幂级数的运算}

设 $\sum a_nx^n=f(x)$，$\sum b_nx^n=g(x)$，收敛半径 $R_1,R_2$，$R=\min\{R_1,R_2\}$，则在 $(-R,R)$ 内
\[
\sum a_nx^n\pm\sum b_nx^n=\sum(a_n\pm b_n)x^n,\qquad
\left(\sum a_nx^n\right)\left(\sum b_nx^n\right)=\sum_{n=0}^{\infty}\left(\sum_{k=0}^{n}a_kb_{n-k}\right)x^n .
\]
在 $(-R,R)$ 内和函数连续，且可\textbf{逐项求导}、\textbf{逐项积分}：
\keyformula{s'(x)=\sum_{n=1}^{\infty}na_nx^{n-1},\qquad \int_0^{x}s(t)\,\mathrm{d}t=\sum_{n=0}^{\infty}\frac{a_n}{n+1}x^{n+1}}
运算后\textbf{收敛半径不变}，但\textbf{端点处的收敛性可能改变}。

\corepoint{逐项求导降次、可能"丢失"常数项，逐项积分升次并添 $x$ 因子——求\textbf{和函数}正靠这两招把系数凑成可积分/可求导的形状，再回到已知展开式。逐项运算只在\textbf{开区间}内保证成立，端点要重新讨论。}

\section{函数展开成幂级数}
\subsection{泰勒级数与展开条件}

若 $f(x)$ 在 $x_0$ 的某邻域内有任意阶导数，则称
\[
\sum_{n=0}^{\infty}\frac{f^{(n)}(x_0)}{n!}(x-x_0)^n=f(x_0)+f'(x_0)(x-x_0)+\frac{f''(x_0)}{2!}(x-x_0)^2+\cdots
\]
为 $f(x)$ 在 $x_0$ 处的\textbf{泰勒级数}；$x_0=0$ 时称\textbf{麦克劳林级数}。

\textbf{定理（展开的充要条件）}\quad $f(x)$ 在该邻域内能展开成泰勒级数 $\iff$ 泰勒公式的余项满足 $\lim\limits_{n\to\infty}R_n(x)=0$。

\textbf{定理（唯一性）}\quad 若 $f(x)$ 能展开成 $x-x_0$ 的幂级数，则系数必为 $a_n=\dfrac{f^{(n)}(x_0)}{n!}$，展开式唯一。

\subsection{五个基本展开式}

\begin{center}
\begin{tabular}{@{}lll@{}}
\toprule
函数 & 展开式 & 收敛域 \\
\midrule
$\mathrm{e}^{x}$ & $\sum\limits_{n=0}^{\infty}\dfrac{x^{n}}{n!}=1+x+\dfrac{x^{2}}{2!}+\cdots$ & $(-\infty,+\infty)$ \\
$\sin x$ & $\sum\limits_{n=0}^{\infty}\dfrac{(-1)^{n}x^{2n+1}}{(2n+1)!}=x-\dfrac{x^{3}}{3!}+\dfrac{x^{5}}{5!}-\cdots$ & $(-\infty,+\infty)$ \\
$\cos x$ & $\sum\limits_{n=0}^{\infty}\dfrac{(-1)^{n}x^{2n}}{(2n)!}=1-\dfrac{x^{2}}{2!}+\dfrac{x^{4}}{4!}-\cdots$ & $(-\infty,+\infty)$ \\
$\ln(1+x)$ & $\sum\limits_{n=1}^{\infty}\dfrac{(-1)^{n-1}x^{n}}{n}=x-\dfrac{x^{2}}{2}+\dfrac{x^{3}}{3}-\cdots$ & $(-1,1]$ \\
$(1+x)^{m}$ & $1+mx+\dfrac{m(m-1)}{2!}x^{2}+\cdots$（$m$ 非非负整数） & $(-1,1)$ \\
\bottomrule
\end{tabular}
\end{center}

\keyformula{\mathrm{e}^{x}=\sum_{n=0}^{\infty}\frac{x^{n}}{n!},\qquad \frac{1}{1-x}=\sum_{n=0}^{\infty}x^{n}\ (|x|<1)}

\subsection{间接展开法}

直接按 $a_n=\frac{f^{(n)}(0)}{n!}$ 求系数（直接展开法）计算量大，通常用\textbf{间接展开法}：由已知展开式经\textbf{变量代换}、四则运算、\textbf{逐项求导}、\textbf{逐项积分}得新展开式，由唯一性保证正确。

\textbf{常用套路}：(1) 把 $f$ 化成 $\mathrm{e}^{u},\sin u,\ln(1+u),(1+u)^m,\frac{1}{1-u}$ 的形式，令 $u=\varphi(x)$；(2) 有理分式先\textbf{拆成部分分式}再用 $\frac{1}{1\mp x}$；(3) 含 $\arctan x,\arcsin x,\ln(1+x)$ 的先求导化成有理式，展开后再逐项积分。

\corepoint{间接展开最易错的是\textbf{收敛域}：代换后范围要跟着换。由 $\frac{1}{1-x}=\sum x^n\,(|x|<1)$ 得 $\frac{1}{1+x^2}=\sum(-1)^nx^{2n}$，收敛域仍是 $|x|<1$；而 $\ln(1+x)$ 逐项积分后在 $x=1$ 也收敛（得 $(-1,1]$）。\textbf{端点是否包含必须重新判}，不能照抄原级数区间。}

\section{函数的幂级数展开式的应用}

利用幂级数可作近似计算，误差由\textbf{余项}估计：(1) \textbf{求函数近似值}——把 $x$ 取小，展开到所需精度，由余项定项数；(2) \textbf{求积分近似值}——被积函数展开后\textbf{逐项积分}再截断；(3) \textbf{求极限}——用展开式代替函数，同阶项相消。

\textbf{微分方程的幂级数解法}：对 $y'=f(x,y)$，若 $f$ 在 $(x_0,y_0)$ 附近可展成幂级数，可设 $y=\sum a_n(x-x_0)^n$，代入比较同次幂系数逐次定出 $a_n$。

\corepoint{核心是\textbf{误差控制}：问"精确到 $10^{-4}$"就要由 $|r_n|\le10^{-4}$ 反解 $n$。交错级数（满足莱布尼茨条件）的余项不超过第一个被舍去项，是最快的估计。}

\section{函数项级数的一致收敛性}

\textbf{定义}\quad 设 $\sum u_n(x)$ 在 $I$ 上收敛于 $s(x)$。若对任意 $\varepsilon>0$，存在\textbf{与 $x$ 无关}的 $N$，使 $n>N$ 时对一切 $x\in I$ 有 $|r_n(x)|<\varepsilon$，则称它在 $I$ 上\textbf{一致收敛}。\textbf{魏尔斯特拉斯判别法}：若 $|u_n(x)|\le M_n$ 且 $\sum M_n$ 收敛，则 $\sum u_n(x)$ 一致收敛。一致收敛级数具有\textbf{连续性}、可\textbf{逐项积分}、可\textbf{逐项求导}（导级数须一致收敛）。

\corepoint{与普通收敛的差别只在 $N$ 是否\textbf{与 $x$ 无关}；有了它才能逐项求极限、积分、求导。幂级数在其收敛区间内部一致收敛，这正是它可逐项运算的依据。}

\section{傅里叶级数}
\subsection{三角级数与三角函数系的正交性}

形如 $\frac{a_0}{2}+\sum\limits_{n=1}^{\infty}(a_n\cos nx+b_n\sin nx)$ 的级数称\textbf{三角级数}。三角函数系 $1,\cos x,\sin x,\cos 2x,\sin 2x,\cdots$ 在 $[-\pi,\pi]$ 上\textbf{正交}：
\[
\int_{-\pi}^{\pi}\cos mx\cos nx\,\mathrm{d}x=\int_{-\pi}^{\pi}\sin mx\sin nx\,\mathrm{d}x=\int_{-\pi}^{\pi}\cos mx\sin nx\,\mathrm{d}x=0\quad(m\ne n),
\]
而 $\int_{-\pi}^{\pi}1^2\mathrm{d}x=2\pi$，$\int_{-\pi}^{\pi}\cos^2nx\,\mathrm{d}x=\int_{-\pi}^{\pi}\sin^2nx\,\mathrm{d}x=\pi$。

\subsection{函数展开成傅里叶级数}

设 $f(x)$ 是周期为 $2\pi$ 的周期函数，由正交性逐项积分得\textbf{欧拉—傅里叶公式}
\keyformula{a_n=\frac{1}{\pi}\int_{-\pi}^{\pi}f(x)\cos nx\,\mathrm{d}x\ (n=0,1,2,\cdots),\qquad b_n=\frac{1}{\pi}\int_{-\pi}^{\pi}f(x)\sin nx\,\mathrm{d}x\ (n=1,2,\cdots)}
由此系数确定的三角级数称 $f$ 的\textbf{傅里叶级数}。

\textbf{狄利克雷收敛定理}\quad 设 $f$ 是周期 $2\pi$ 的周期函数，若 (1) 一个周期内\textbf{连续或只有有限个第一类间断点}；(2) 一个周期内\textbf{至多有有限个极值点}，则其傅里叶级数收敛，且
\[
\frac{a_0}{2}+\sum_{n=1}^{\infty}(a_n\cos nx+b_n\sin nx)=
\begin{cases}
f(x), & x\ \text{为连续点},\\[3pt]
\dfrac{f(x^-)+f(x^+)}{2}, & x\ \text{为间断点}.
\end{cases}
\]

\corepoint{两个要点：(1) 系数公式的 $\frac{1}{\pi}$ 与 $a_0$ 前的 $\frac12$ 配套，写错全错；(2) 收敛定理给的是"和函数"——\textbf{间断点}处收敛于左右极限的\textbf{平均值}，端点 $x=\pm\pi$ 处收敛于 $\frac{f(-\pi^+)+f(\pi^-)}{2}$。问"展开式在某点等于什么"，往往就考这个平均值。}

\subsection{正弦级数和余弦级数}

\begin{itemize}
  \item $f$ 为\textbf{奇函数}：$a_n=0$，$b_n=\frac{2}{\pi}\int_0^{\pi}f(x)\sin nx\,\mathrm{d}x$，级数只含正弦项（\textbf{正弦级数}）。
  \item $f$ 为\textbf{偶函数}：$b_n=0$，$a_n=\frac{2}{\pi}\int_0^{\pi}f(x)\cos nx\,\mathrm{d}x$，级数只含余弦项（\textbf{余弦级数}）。
\end{itemize}

\textbf{有限区间上的展开（奇延拓与偶延拓）}：若 $f$ 只在 $[0,\pi]$ 上有定义，要展成正弦级数就作\textbf{奇延拓}（$F(-x)=-F(x)$）再按奇函数求 $b_n$；要展成余弦级数就作\textbf{偶延拓}（$F(-x)=F(x)$）再按偶函数求 $a_n$。所得级数在 $(0,\pi)$ 内收敛于 $f(x)$。

\corepoint{奇延拓的正弦级数在 $x=0$ 处收敛于 $0$（除非 $f(0)=0$）；偶延拓的余弦级数在 $x=0$ 处收敛于 $f(0^+)$。这是"端点值"类题目的标准考法；延拓只是求系数的技巧，原函数在 $[0,\pi]$ 上并未改变。}

\section{一般周期函数的傅里叶级数}
\subsection{周期为 $2l$ 的周期函数}

设 $f(x)$ 的周期为 $2l$，令 $t=\frac{\pi x}{l}$ 可化为周期 $2\pi$ 的情形，得
\keyformula{f(x)\sim\frac{a_0}{2}+\sum_{n=1}^{\infty}\left(a_n\cos\frac{n\pi x}{l}+b_n\sin\frac{n\pi x}{l}\right)}
其中
\[
a_n=\frac{1}{l}\int_{-l}^{l}f(x)\cos\frac{n\pi x}{l}\,\mathrm{d}x\ (n=0,1,2,\cdots),\qquad
b_n=\frac{1}{l}\int_{-l}^{l}f(x)\sin\frac{n\pi x}{l}\,\mathrm{d}x\ (n=1,2,\cdots).
\]
$f$ 为奇函数时 $a_n=0$，$b_n=\frac{2}{l}\int_0^{l}f\sin\frac{n\pi x}{l}\mathrm{d}x$；为偶函数时 $b_n=0$，$a_n=\frac{2}{l}\int_0^{l}f\cos\frac{n\pi x}{l}\mathrm{d}x$。狄利克雷定理同样成立。

\subsection{傅里叶级数的复数形式（简）}

由欧拉公式 $\mathrm{e}^{\mathrm{i}\theta}=\cos\theta+\mathrm{i}\sin\theta$ 得复数形式
\keyformula{f(x)\sim\sum_{n=-\infty}^{+\infty}c_n\mathrm{e}^{\mathrm{i}\frac{n\pi x}{l}},\qquad c_n=\frac{1}{2l}\int_{-l}^{l}f(x)\,\mathrm{e}^{-\mathrm{i}\frac{n\pi x}{l}}\,\mathrm{d}x\ (n=0,\pm1,\pm2,\cdots)}
周期为 $2\pi$ 时 $c_n=\frac{1}{2\pi}\int_{-\pi}^{\pi}f(x)\mathrm{e}^{-\mathrm{i}nx}\mathrm{d}x$，且 $c_0=\frac{a_0}{2}$，$c_n=\frac{a_n-\mathrm{i}b_n}{2}$，$c_{-n}=\frac{a_n+\mathrm{i}b_n}{2}$。

\corepoint{周期 $2l$ 的记法：把 $2\pi$ 情形的 $\cos nx$ 换成 $\cos\frac{n\pi x}{l}$，系数前的 $\frac{1}{\pi}$ 换成 $\frac{1}{l}$（区间长度由 $2\pi$ 变 $2l$）。\textbf{积分区间须与周期一致}（$[-l,l]$ 或任一长 $2l$ 的区间），否则系数全错。}
 引入，不含 document 环境

\section{常数项级数的概念和性质}
\subsection{概念}

\textbf{定义}\quad 给定数列 $u_1,u_2,\cdots,u_n,\cdots$，称 $\sum\limits_{n=1}^{\infty}u_n=u_1+u_2+\cdots+u_n+\cdots$ 为（常数项）\textbf{无穷级数}，$u_n$ 为一般项。作部分和 $s_n=u_1+\cdots+u_n$，若 $\lim\limits_{n\to\infty}s_n=s$（常数），称级数\textbf{收敛}，$s$ 为\textbf{和}；若 $\{s_n\}$ 无极限，称级数\textbf{发散}。收敛时 $r_n=s-s_n=\sum\limits_{k=n+1}^{\infty}u_k$ 称\textbf{余项}，且 $r_n\to0$。

\textbf{两个基准级数}（必须记牢）：
\keyformula{\sum_{n=0}^{\infty}aq^{n}=\frac{a}{1-q}\ (|q|<1),\quad \text{发散}\ (|q|\ge1);\qquad \sum_{n=1}^{\infty}\frac{1}{n^{p}}\ \text{收敛}\iff p>1}
几何级数：$|q|<1$ 收敛于 $\frac{a}{1-q}$；$p$ 级数：$p>1$ 收敛，$p\le1$ 发散（$p=1$ 为调和级数，发散）。

\subsection{收敛级数的基本性质}

\begin{enumerate}
  \item \textbf{线性}：$\sum u_n=s,\ \sum v_n=t\Rightarrow\sum(\alpha u_n+\beta v_n)=\alpha s+\beta t$。
  \item \textbf{去掉、加上或改变有限项}不改变收敛性（收敛时和一般会变）。
  \item \textbf{收敛的必要条件}：若 $\sum u_n$ 收敛，则 $\lim\limits_{n\to\infty}u_n=0$。
  \item \textbf{收敛级数可任意加括号}（和不变）；但\textbf{去括号}不一定成立。
  \item \textbf{柯西审敛原理}：$\sum u_n$ 收敛 $\iff$ 对任意 $\varepsilon>0$ 存在 $N$，当 $n>N$ 时对任意正整数 $p$ 有 $|u_{n+1}+\cdots+u_{n+p}|<\varepsilon$。
\end{enumerate}

\corepoint{性质 3 是\textbf{必要}条件，只用来判\textbf{发散}：$\lim u_n\ne0$ 则必发散；但 $\lim u_n=0$ 推不出收敛（调和级数 $\sum\frac1n$ 即反例），这是判敛第一步。性质 4 的单向性常考：加括号收敛 $\not\Rightarrow$ 原级数收敛。}

\section{常数项级数的审敛法}
\subsection{正项级数及其审敛法}

若 $u_n\ge0$，称 $\sum u_n$ 为\textbf{正项级数}。

\textbf{定理 1（基本定理）}\quad 正项级数收敛的\textbf{充要条件}是部分和数列 $\{s_n\}$ \textbf{有界}（$s_n$ 单调增，单调有界必收敛）。

\textbf{定理 2（比较审敛法）}\quad 设 $\sum u_n,\sum v_n$ 为正项级数且 $u_n\le v_n$：若 $\sum v_n$ 收敛则 $\sum u_n$ 收敛；若 $\sum u_n$ 发散则 $\sum v_n$ 发散。

\textbf{推论（极限形式）}\quad 设 $\lim\limits_{n\to\infty}\frac{u_n}{v_n}=l$：$0<l<+\infty$ 时两者\textbf{同敛散}；$l=0$ 且 $\sum v_n$ 收敛则 $\sum u_n$ 收敛；$l=+\infty$ 且 $\sum v_n$ 发散则 $\sum u_n$ 发散。

\textbf{定理 3（比值审敛法，达朗贝尔）}\quad 对正项级数 $\sum u_n$，若 $\lim\limits_{n\to\infty}\frac{u_{n+1}}{u_n}=\rho$，则 $\rho<1$ 收敛，$\rho>1$（含 $+\infty$）发散，$\rho=1$ \textbf{失效}。

\textbf{定理 4（根值审敛法，柯西）}\quad 若 $\lim\limits_{n\to\infty}\sqrt[n]{u_n}=\rho$，结论同定理 3。

\textbf{定理 5（积分审敛法）}\quad 设 $f(x)$ 在 $[1,+\infty)$ 上连续、非负、单调减少且 $f(n)=u_n$，则 $\sum\limits_{n=1}^{\infty}u_n$ 与 $\int_1^{+\infty}f(x)\,\mathrm{d}x$ \textbf{同时收敛或同时发散}。

\corepoint{比较法的\textbf{关键是选比较对象}：常用几何级数、$p$ 级数、$\frac{1}{n\ln^pn}$ 型。选法看通项的阶：含 $n!$、$a^n$、$n^n$ 优先用\textbf{比值}；通项为 $(f(n))^n$ 形优先用\textbf{根值}；$n$ 的有理分式（含根号）优先用\textbf{比较法}（与 $\frac{1}{n^p}$ 比，看分母最高次与分子之差）。比值、根值法在 $\rho=1$ 时\textbf{不能用}（$\sum\frac1n$、$\sum\frac{1}{n^2}$ 都给出 $\rho=1$），这是最常见失分点。}

\subsection{交错级数及其审敛法}

形如 $\sum\limits_{n=1}^{\infty}(-1)^{n-1}u_n=u_1-u_2+u_3-\cdots$（$u_n>0$）的级数称\textbf{交错级数}。

\textbf{定理 6（莱布尼茨判别法）}\quad 若 (1) $u_n\ge u_{n+1}$（$n=1,2,\cdots$）；(2) $\lim\limits_{n\to\infty}u_n=0$，则 $\sum(-1)^{n-1}u_n$ \textbf{收敛}，且 $s\le u_1$，余项 $|r_n|\le u_{n+1}$。

\corepoint{两条件缺一不可：$u_n$ \textbf{单调不增}且 $\to0$；只看 $\to0$ 不够（不单调时可能发散）。非交错级数不能套用，应先拆项。余项估计 $|r_n|\le u_{n+1}$ 常用于近似计算定误差。}

\subsection{绝对收敛与条件收敛}

对任意项级数 $\sum u_n$：若 $\sum|u_n|$ 收敛，称 $\sum u_n$ \textbf{绝对收敛}；若 $\sum u_n$ 收敛而 $\sum|u_n|$ 发散，称\textbf{条件收敛}。

\textbf{定理 7}\quad 绝对收敛 $\Rightarrow$ 收敛。

\textbf{定理 8（绝对收敛级数的性质）}\quad 绝对收敛级数可\textbf{任意交换项次序}而不改变其和；两个绝对收敛级数的\textbf{柯西乘积} $\sum\limits_{n=1}^{\infty}\left(\sum\limits_{k=1}^{n}u_kv_{n-k+1}\right)$ 也绝对收敛，和为 $(\sum u_n)(\sum v_n)$。条件收敛级数不能随意重排。

\corepoint{判任意项级数的\textbf{标准流程}：先看 $\sum|u_n|$——收敛即\textbf{绝对收敛}，结束；若发散，再判原级数（交错用莱布尼茨，一般用定义）——收敛则为\textbf{条件收敛}。$|u_n|$ 发散\textbf{不能}推出 $u_n$ 发散，这是最易搞反的一步。}

\section{幂级数}
\subsection{函数项级数}

设 $u_n(x)$ 定义在区间 $I$ 上，称 $\sum\limits_{n=1}^{\infty}u_n(x)$ 为\textbf{函数项级数}。若 $\sum u_n(x_0)$ 收敛，称 $x_0$ 为收敛点，收敛点全体称\textbf{收敛域}，其上的 $s(x)=\sum u_n(x)$ 称\textbf{和函数}。

\subsection{幂级数及其收敛性}

形如 $\sum\limits_{n=0}^{\infty}a_n(x-x_0)^n$ 的级数称\textbf{幂级数}（$x_0=0$ 时即 $\sum a_nx^n$）。

\textbf{阿贝尔定理}\quad 若 $\sum a_nx^n$ 在 $x=x_0\ne0$ 处收敛，则对 $|x|<|x_0|$ 的一切 $x$ \textbf{绝对收敛}；若在 $x_0$ 处发散，则对 $|x|>|x_0|$ 的一切 $x$ \textbf{发散}。

故收敛性只有三种：仅在 $x=0$ 收敛（$R=0$）；全数轴收敛（$R=+\infty$）；存在 $R>0$，$|x|<R$ 绝对收敛、$|x|>R$ 发散——$R$ 称\textbf{收敛半径}，$(-R,R)$ 称\textbf{收敛区间}，端点 $x=\pm R$ 须单独判定，最终得\textbf{收敛域}。

\textbf{收敛半径的求法}：对 $\sum a_nx^n$，若
\keyformula{\rho=\lim_{n\to\infty}\left|\frac{a_{n+1}}{a_n}\right|\quad\text{或}\quad \rho=\lim_{n\to\infty}\sqrt[n]{|a_n|},\qquad R=\frac{1}{\rho}}
（约定 $\rho=0$ 时 $R=+\infty$；$\rho=+\infty$ 时 $R=0$。）

\corepoint{求收敛域\textbf{三步}：(1) 用公式求 $R$（$a_n$ 是 $x^n$ 的系数；缺项级数如 $\sum a_nx^{2n}$ 不能直接套公式，应把 $x^2$ 当整体用比值法）；(2) 得 $(-R,R)$；(3) \textbf{端点单独讨论}——代入 $x=\pm R$ 化成数项级数再判，两端可能一敛一散。只写开区间不讨论端点是最常见的扣分点。}

\subsection{幂级数的运算}

设 $\sum a_nx^n=f(x)$，$\sum b_nx^n=g(x)$，收敛半径 $R_1,R_2$，$R=\min\{R_1,R_2\}$，则在 $(-R,R)$ 内
\[
\sum a_nx^n\pm\sum b_nx^n=\sum(a_n\pm b_n)x^n,\qquad
\left(\sum a_nx^n\right)\left(\sum b_nx^n\right)=\sum_{n=0}^{\infty}\left(\sum_{k=0}^{n}a_kb_{n-k}\right)x^n .
\]
在 $(-R,R)$ 内和函数连续，且可\textbf{逐项求导}、\textbf{逐项积分}：
\keyformula{s'(x)=\sum_{n=1}^{\infty}na_nx^{n-1},\qquad \int_0^{x}s(t)\,\mathrm{d}t=\sum_{n=0}^{\infty}\frac{a_n}{n+1}x^{n+1}}
运算后\textbf{收敛半径不变}，但\textbf{端点处的收敛性可能改变}。

\corepoint{逐项求导降次、可能"丢失"常数项，逐项积分升次并添 $x$ 因子——求\textbf{和函数}正靠这两招把系数凑成可积分/可求导的形状，再回到已知展开式。逐项运算只在\textbf{开区间}内保证成立，端点要重新讨论。}

\section{函数展开成幂级数}
\subsection{泰勒级数与展开条件}

若 $f(x)$ 在 $x_0$ 的某邻域内有任意阶导数，则称
\[
\sum_{n=0}^{\infty}\frac{f^{(n)}(x_0)}{n!}(x-x_0)^n=f(x_0)+f'(x_0)(x-x_0)+\frac{f''(x_0)}{2!}(x-x_0)^2+\cdots
\]
为 $f(x)$ 在 $x_0$ 处的\textbf{泰勒级数}；$x_0=0$ 时称\textbf{麦克劳林级数}。

\textbf{定理（展开的充要条件）}\quad $f(x)$ 在该邻域内能展开成泰勒级数 $\iff$ 泰勒公式的余项满足 $\lim\limits_{n\to\infty}R_n(x)=0$。

\textbf{定理（唯一性）}\quad 若 $f(x)$ 能展开成 $x-x_0$ 的幂级数，则系数必为 $a_n=\dfrac{f^{(n)}(x_0)}{n!}$，展开式唯一。

\subsection{五个基本展开式}

\begin{center}
\begin{tabular}{@{}lll@{}}
\toprule
函数 & 展开式 & 收敛域 \\
\midrule
$\mathrm{e}^{x}$ & $\sum\limits_{n=0}^{\infty}\dfrac{x^{n}}{n!}=1+x+\dfrac{x^{2}}{2!}+\cdots$ & $(-\infty,+\infty)$ \\
$\sin x$ & $\sum\limits_{n=0}^{\infty}\dfrac{(-1)^{n}x^{2n+1}}{(2n+1)!}=x-\dfrac{x^{3}}{3!}+\dfrac{x^{5}}{5!}-\cdots$ & $(-\infty,+\infty)$ \\
$\cos x$ & $\sum\limits_{n=0}^{\infty}\dfrac{(-1)^{n}x^{2n}}{(2n)!}=1-\dfrac{x^{2}}{2!}+\dfrac{x^{4}}{4!}-\cdots$ & $(-\infty,+\infty)$ \\
$\ln(1+x)$ & $\sum\limits_{n=1}^{\infty}\dfrac{(-1)^{n-1}x^{n}}{n}=x-\dfrac{x^{2}}{2}+\dfrac{x^{3}}{3}-\cdots$ & $(-1,1]$ \\
$(1+x)^{m}$ & $1+mx+\dfrac{m(m-1)}{2!}x^{2}+\cdots$（$m$ 非非负整数） & $(-1,1)$ \\
\bottomrule
\end{tabular}
\end{center}

\keyformula{\mathrm{e}^{x}=\sum_{n=0}^{\infty}\frac{x^{n}}{n!},\qquad \frac{1}{1-x}=\sum_{n=0}^{\infty}x^{n}\ (|x|<1)}

\subsection{间接展开法}

直接按 $a_n=\frac{f^{(n)}(0)}{n!}$ 求系数（直接展开法）计算量大，通常用\textbf{间接展开法}：由已知展开式经\textbf{变量代换}、四则运算、\textbf{逐项求导}、\textbf{逐项积分}得新展开式，由唯一性保证正确。

\textbf{常用套路}：(1) 把 $f$ 化成 $\mathrm{e}^{u},\sin u,\ln(1+u),(1+u)^m,\frac{1}{1-u}$ 的形式，令 $u=\varphi(x)$；(2) 有理分式先\textbf{拆成部分分式}再用 $\frac{1}{1\mp x}$；(3) 含 $\arctan x,\arcsin x,\ln(1+x)$ 的先求导化成有理式，展开后再逐项积分。

\corepoint{间接展开最易错的是\textbf{收敛域}：代换后范围要跟着换。由 $\frac{1}{1-x}=\sum x^n\,(|x|<1)$ 得 $\frac{1}{1+x^2}=\sum(-1)^nx^{2n}$，收敛域仍是 $|x|<1$；而 $\ln(1+x)$ 逐项积分后在 $x=1$ 也收敛（得 $(-1,1]$）。\textbf{端点是否包含必须重新判}，不能照抄原级数区间。}

\section{函数的幂级数展开式的应用}

利用幂级数可作近似计算，误差由\textbf{余项}估计：(1) \textbf{求函数近似值}——把 $x$ 取小，展开到所需精度，由余项定项数；(2) \textbf{求积分近似值}——被积函数展开后\textbf{逐项积分}再截断；(3) \textbf{求极限}——用展开式代替函数，同阶项相消。

\textbf{微分方程的幂级数解法}：对 $y'=f(x,y)$，若 $f$ 在 $(x_0,y_0)$ 附近可展成幂级数，可设 $y=\sum a_n(x-x_0)^n$，代入比较同次幂系数逐次定出 $a_n$。

\corepoint{核心是\textbf{误差控制}：问"精确到 $10^{-4}$"就要由 $|r_n|\le10^{-4}$ 反解 $n$。交错级数（满足莱布尼茨条件）的余项不超过第一个被舍去项，是最快的估计。}

\section{函数项级数的一致收敛性}

\textbf{定义}\quad 设 $\sum u_n(x)$ 在 $I$ 上收敛于 $s(x)$。若对任意 $\varepsilon>0$，存在\textbf{与 $x$ 无关}的 $N$，使 $n>N$ 时对一切 $x\in I$ 有 $|r_n(x)|<\varepsilon$，则称它在 $I$ 上\textbf{一致收敛}。\textbf{魏尔斯特拉斯判别法}：若 $|u_n(x)|\le M_n$ 且 $\sum M_n$ 收敛，则 $\sum u_n(x)$ 一致收敛。一致收敛级数具有\textbf{连续性}、可\textbf{逐项积分}、可\textbf{逐项求导}（导级数须一致收敛）。

\corepoint{与普通收敛的差别只在 $N$ 是否\textbf{与 $x$ 无关}；有了它才能逐项求极限、积分、求导。幂级数在其收敛区间内部一致收敛，这正是它可逐项运算的依据。}

\section{傅里叶级数}
\subsection{三角级数与三角函数系的正交性}

形如 $\frac{a_0}{2}+\sum\limits_{n=1}^{\infty}(a_n\cos nx+b_n\sin nx)$ 的级数称\textbf{三角级数}。三角函数系 $1,\cos x,\sin x,\cos 2x,\sin 2x,\cdots$ 在 $[-\pi,\pi]$ 上\textbf{正交}：
\[
\int_{-\pi}^{\pi}\cos mx\cos nx\,\mathrm{d}x=\int_{-\pi}^{\pi}\sin mx\sin nx\,\mathrm{d}x=\int_{-\pi}^{\pi}\cos mx\sin nx\,\mathrm{d}x=0\quad(m\ne n),
\]
而 $\int_{-\pi}^{\pi}1^2\mathrm{d}x=2\pi$，$\int_{-\pi}^{\pi}\cos^2nx\,\mathrm{d}x=\int_{-\pi}^{\pi}\sin^2nx\,\mathrm{d}x=\pi$。

\subsection{函数展开成傅里叶级数}

设 $f(x)$ 是周期为 $2\pi$ 的周期函数，由正交性逐项积分得\textbf{欧拉—傅里叶公式}
\keyformula{a_n=\frac{1}{\pi}\int_{-\pi}^{\pi}f(x)\cos nx\,\mathrm{d}x\ (n=0,1,2,\cdots),\qquad b_n=\frac{1}{\pi}\int_{-\pi}^{\pi}f(x)\sin nx\,\mathrm{d}x\ (n=1,2,\cdots)}
由此系数确定的三角级数称 $f$ 的\textbf{傅里叶级数}。

\textbf{狄利克雷收敛定理}\quad 设 $f$ 是周期 $2\pi$ 的周期函数，若 (1) 一个周期内\textbf{连续或只有有限个第一类间断点}；(2) 一个周期内\textbf{至多有有限个极值点}，则其傅里叶级数收敛，且
\[
\frac{a_0}{2}+\sum_{n=1}^{\infty}(a_n\cos nx+b_n\sin nx)=
\begin{cases}
f(x), & x\ \text{为连续点},\\[3pt]
\dfrac{f(x^-)+f(x^+)}{2}, & x\ \text{为间断点}.
\end{cases}
\]

\corepoint{两个要点：(1) 系数公式的 $\frac{1}{\pi}$ 与 $a_0$ 前的 $\frac12$ 配套，写错全错；(2) 收敛定理给的是"和函数"——\textbf{间断点}处收敛于左右极限的\textbf{平均值}，端点 $x=\pm\pi$ 处收敛于 $\frac{f(-\pi^+)+f(\pi^-)}{2}$。问"展开式在某点等于什么"，往往就考这个平均值。}

\subsection{正弦级数和余弦级数}

\begin{itemize}
  \item $f$ 为\textbf{奇函数}：$a_n=0$，$b_n=\frac{2}{\pi}\int_0^{\pi}f(x)\sin nx\,\mathrm{d}x$，级数只含正弦项（\textbf{正弦级数}）。
  \item $f$ 为\textbf{偶函数}：$b_n=0$，$a_n=\frac{2}{\pi}\int_0^{\pi}f(x)\cos nx\,\mathrm{d}x$，级数只含余弦项（\textbf{余弦级数}）。
\end{itemize}

\textbf{有限区间上的展开（奇延拓与偶延拓）}：若 $f$ 只在 $[0,\pi]$ 上有定义，要展成正弦级数就作\textbf{奇延拓}（$F(-x)=-F(x)$）再按奇函数求 $b_n$；要展成余弦级数就作\textbf{偶延拓}（$F(-x)=F(x)$）再按偶函数求 $a_n$。所得级数在 $(0,\pi)$ 内收敛于 $f(x)$。

\corepoint{奇延拓的正弦级数在 $x=0$ 处收敛于 $0$（除非 $f(0)=0$）；偶延拓的余弦级数在 $x=0$ 处收敛于 $f(0^+)$。这是"端点值"类题目的标准考法；延拓只是求系数的技巧，原函数在 $[0,\pi]$ 上并未改变。}

\section{一般周期函数的傅里叶级数}
\subsection{周期为 $2l$ 的周期函数}

设 $f(x)$ 的周期为 $2l$，令 $t=\frac{\pi x}{l}$ 可化为周期 $2\pi$ 的情形，得
\keyformula{f(x)\sim\frac{a_0}{2}+\sum_{n=1}^{\infty}\left(a_n\cos\frac{n\pi x}{l}+b_n\sin\frac{n\pi x}{l}\right)}
其中
\[
a_n=\frac{1}{l}\int_{-l}^{l}f(x)\cos\frac{n\pi x}{l}\,\mathrm{d}x\ (n=0,1,2,\cdots),\qquad
b_n=\frac{1}{l}\int_{-l}^{l}f(x)\sin\frac{n\pi x}{l}\,\mathrm{d}x\ (n=1,2,\cdots).
\]
$f$ 为奇函数时 $a_n=0$，$b_n=\frac{2}{l}\int_0^{l}f\sin\frac{n\pi x}{l}\mathrm{d}x$；为偶函数时 $b_n=0$，$a_n=\frac{2}{l}\int_0^{l}f\cos\frac{n\pi x}{l}\mathrm{d}x$。狄利克雷定理同样成立。

\subsection{傅里叶级数的复数形式（简）}

由欧拉公式 $\mathrm{e}^{\mathrm{i}\theta}=\cos\theta+\mathrm{i}\sin\theta$ 得复数形式
\keyformula{f(x)\sim\sum_{n=-\infty}^{+\infty}c_n\mathrm{e}^{\mathrm{i}\frac{n\pi x}{l}},\qquad c_n=\frac{1}{2l}\int_{-l}^{l}f(x)\,\mathrm{e}^{-\mathrm{i}\frac{n\pi x}{l}}\,\mathrm{d}x\ (n=0,\pm1,\pm2,\cdots)}
周期为 $2\pi$ 时 $c_n=\frac{1}{2\pi}\int_{-\pi}^{\pi}f(x)\mathrm{e}^{-\mathrm{i}nx}\mathrm{d}x$，且 $c_0=\frac{a_0}{2}$，$c_n=\frac{a_n-\mathrm{i}b_n}{2}$，$c_{-n}=\frac{a_n+\mathrm{i}b_n}{2}$。

\corepoint{周期 $2l$ 的记法：把 $2\pi$ 情形的 $\cos nx$ 换成 $\cos\frac{n\pi x}{l}$，系数前的 $\frac{1}{\pi}$ 换成 $\frac{1}{l}$（区间长度由 $2\pi$ 变 $2l$）。\textbf{积分区间须与周期一致}（$[-l,l]$ 或任一长 $2l$ 的区间），否则系数全错。}
 引入，不含 document 环境

\section{常数项级数的概念和性质}
\subsection{概念}

\textbf{定义}\quad 给定数列 $u_1,u_2,\cdots,u_n,\cdots$，称 $\sum\limits_{n=1}^{\infty}u_n=u_1+u_2+\cdots+u_n+\cdots$ 为（常数项）\textbf{无穷级数}，$u_n$ 为一般项。作部分和 $s_n=u_1+\cdots+u_n$，若 $\lim\limits_{n\to\infty}s_n=s$（常数），称级数\textbf{收敛}，$s$ 为\textbf{和}；若 $\{s_n\}$ 无极限，称级数\textbf{发散}。收敛时 $r_n=s-s_n=\sum\limits_{k=n+1}^{\infty}u_k$ 称\textbf{余项}，且 $r_n\to0$。

\textbf{两个基准级数}（必须记牢）：
\keyformula{\sum_{n=0}^{\infty}aq^{n}=\frac{a}{1-q}\ (|q|<1),\quad \text{发散}\ (|q|\ge1);\qquad \sum_{n=1}^{\infty}\frac{1}{n^{p}}\ \text{收敛}\iff p>1}
几何级数：$|q|<1$ 收敛于 $\frac{a}{1-q}$；$p$ 级数：$p>1$ 收敛，$p\le1$ 发散（$p=1$ 为调和级数，发散）。

\subsection{收敛级数的基本性质}

\begin{enumerate}
  \item \textbf{线性}：$\sum u_n=s,\ \sum v_n=t\Rightarrow\sum(\alpha u_n+\beta v_n)=\alpha s+\beta t$。
  \item \textbf{去掉、加上或改变有限项}不改变收敛性（收敛时和一般会变）。
  \item \textbf{收敛的必要条件}：若 $\sum u_n$ 收敛，则 $\lim\limits_{n\to\infty}u_n=0$。
  \item \textbf{收敛级数可任意加括号}（和不变）；但\textbf{去括号}不一定成立。
  \item \textbf{柯西审敛原理}：$\sum u_n$ 收敛 $\iff$ 对任意 $\varepsilon>0$ 存在 $N$，当 $n>N$ 时对任意正整数 $p$ 有 $|u_{n+1}+\cdots+u_{n+p}|<\varepsilon$。
\end{enumerate}

\corepoint{性质 3 是\textbf{必要}条件，只用来判\textbf{发散}：$\lim u_n\ne0$ 则必发散；但 $\lim u_n=0$ 推不出收敛（调和级数 $\sum\frac1n$ 即反例），这是判敛第一步。性质 4 的单向性常考：加括号收敛 $\not\Rightarrow$ 原级数收敛。}

\section{常数项级数的审敛法}
\subsection{正项级数及其审敛法}

若 $u_n\ge0$，称 $\sum u_n$ 为\textbf{正项级数}。

\textbf{定理 1（基本定理）}\quad 正项级数收敛的\textbf{充要条件}是部分和数列 $\{s_n\}$ \textbf{有界}（$s_n$ 单调增，单调有界必收敛）。

\textbf{定理 2（比较审敛法）}\quad 设 $\sum u_n,\sum v_n$ 为正项级数且 $u_n\le v_n$：若 $\sum v_n$ 收敛则 $\sum u_n$ 收敛；若 $\sum u_n$ 发散则 $\sum v_n$ 发散。

\textbf{推论（极限形式）}\quad 设 $\lim\limits_{n\to\infty}\frac{u_n}{v_n}=l$：$0<l<+\infty$ 时两者\textbf{同敛散}；$l=0$ 且 $\sum v_n$ 收敛则 $\sum u_n$ 收敛；$l=+\infty$ 且 $\sum v_n$ 发散则 $\sum u_n$ 发散。

\textbf{定理 3（比值审敛法，达朗贝尔）}\quad 对正项级数 $\sum u_n$，若 $\lim\limits_{n\to\infty}\frac{u_{n+1}}{u_n}=\rho$，则 $\rho<1$ 收敛，$\rho>1$（含 $+\infty$）发散，$\rho=1$ \textbf{失效}。

\textbf{定理 4（根值审敛法，柯西）}\quad 若 $\lim\limits_{n\to\infty}\sqrt[n]{u_n}=\rho$，结论同定理 3。

\textbf{定理 5（积分审敛法）}\quad 设 $f(x)$ 在 $[1,+\infty)$ 上连续、非负、单调减少且 $f(n)=u_n$，则 $\sum\limits_{n=1}^{\infty}u_n$ 与 $\int_1^{+\infty}f(x)\,\mathrm{d}x$ \textbf{同时收敛或同时发散}。

\corepoint{比较法的\textbf{关键是选比较对象}：常用几何级数、$p$ 级数、$\frac{1}{n\ln^pn}$ 型。选法看通项的阶：含 $n!$、$a^n$、$n^n$ 优先用\textbf{比值}；通项为 $(f(n))^n$ 形优先用\textbf{根值}；$n$ 的有理分式（含根号）优先用\textbf{比较法}（与 $\frac{1}{n^p}$ 比，看分母最高次与分子之差）。比值、根值法在 $\rho=1$ 时\textbf{不能用}（$\sum\frac1n$、$\sum\frac{1}{n^2}$ 都给出 $\rho=1$），这是最常见失分点。}

\subsection{交错级数及其审敛法}

形如 $\sum\limits_{n=1}^{\infty}(-1)^{n-1}u_n=u_1-u_2+u_3-\cdots$（$u_n>0$）的级数称\textbf{交错级数}。

\textbf{定理 6（莱布尼茨判别法）}\quad 若 (1) $u_n\ge u_{n+1}$（$n=1,2,\cdots$）；(2) $\lim\limits_{n\to\infty}u_n=0$，则 $\sum(-1)^{n-1}u_n$ \textbf{收敛}，且 $s\le u_1$，余项 $|r_n|\le u_{n+1}$。

\corepoint{两条件缺一不可：$u_n$ \textbf{单调不增}且 $\to0$；只看 $\to0$ 不够（不单调时可能发散）。非交错级数不能套用，应先拆项。余项估计 $|r_n|\le u_{n+1}$ 常用于近似计算定误差。}

\subsection{绝对收敛与条件收敛}

对任意项级数 $\sum u_n$：若 $\sum|u_n|$ 收敛，称 $\sum u_n$ \textbf{绝对收敛}；若 $\sum u_n$ 收敛而 $\sum|u_n|$ 发散，称\textbf{条件收敛}。

\textbf{定理 7}\quad 绝对收敛 $\Rightarrow$ 收敛。

\textbf{定理 8（绝对收敛级数的性质）}\quad 绝对收敛级数可\textbf{任意交换项次序}而不改变其和；两个绝对收敛级数的\textbf{柯西乘积} $\sum\limits_{n=1}^{\infty}\left(\sum\limits_{k=1}^{n}u_kv_{n-k+1}\right)$ 也绝对收敛，和为 $(\sum u_n)(\sum v_n)$。条件收敛级数不能随意重排。

\corepoint{判任意项级数的\textbf{标准流程}：先看 $\sum|u_n|$——收敛即\textbf{绝对收敛}，结束；若发散，再判原级数（交错用莱布尼茨，一般用定义）——收敛则为\textbf{条件收敛}。$|u_n|$ 发散\textbf{不能}推出 $u_n$ 发散，这是最易搞反的一步。}

\section{幂级数}
\subsection{函数项级数}

设 $u_n(x)$ 定义在区间 $I$ 上，称 $\sum\limits_{n=1}^{\infty}u_n(x)$ 为\textbf{函数项级数}。若 $\sum u_n(x_0)$ 收敛，称 $x_0$ 为收敛点，收敛点全体称\textbf{收敛域}，其上的 $s(x)=\sum u_n(x)$ 称\textbf{和函数}。

\subsection{幂级数及其收敛性}

形如 $\sum\limits_{n=0}^{\infty}a_n(x-x_0)^n$ 的级数称\textbf{幂级数}（$x_0=0$ 时即 $\sum a_nx^n$）。

\textbf{阿贝尔定理}\quad 若 $\sum a_nx^n$ 在 $x=x_0\ne0$ 处收敛，则对 $|x|<|x_0|$ 的一切 $x$ \textbf{绝对收敛}；若在 $x_0$ 处发散，则对 $|x|>|x_0|$ 的一切 $x$ \textbf{发散}。

故收敛性只有三种：仅在 $x=0$ 收敛（$R=0$）；全数轴收敛（$R=+\infty$）；存在 $R>0$，$|x|<R$ 绝对收敛、$|x|>R$ 发散——$R$ 称\textbf{收敛半径}，$(-R,R)$ 称\textbf{收敛区间}，端点 $x=\pm R$ 须单独判定，最终得\textbf{收敛域}。

\textbf{收敛半径的求法}：对 $\sum a_nx^n$，若
\keyformula{\rho=\lim_{n\to\infty}\left|\frac{a_{n+1}}{a_n}\right|\quad\text{或}\quad \rho=\lim_{n\to\infty}\sqrt[n]{|a_n|},\qquad R=\frac{1}{\rho}}
（约定 $\rho=0$ 时 $R=+\infty$；$\rho=+\infty$ 时 $R=0$。）

\corepoint{求收敛域\textbf{三步}：(1) 用公式求 $R$（$a_n$ 是 $x^n$ 的系数；缺项级数如 $\sum a_nx^{2n}$ 不能直接套公式，应把 $x^2$ 当整体用比值法）；(2) 得 $(-R,R)$；(3) \textbf{端点单独讨论}——代入 $x=\pm R$ 化成数项级数再判，两端可能一敛一散。只写开区间不讨论端点是最常见的扣分点。}

\subsection{幂级数的运算}

设 $\sum a_nx^n=f(x)$，$\sum b_nx^n=g(x)$，收敛半径 $R_1,R_2$，$R=\min\{R_1,R_2\}$，则在 $(-R,R)$ 内
\[
\sum a_nx^n\pm\sum b_nx^n=\sum(a_n\pm b_n)x^n,\qquad
\left(\sum a_nx^n\right)\left(\sum b_nx^n\right)=\sum_{n=0}^{\infty}\left(\sum_{k=0}^{n}a_kb_{n-k}\right)x^n .
\]
在 $(-R,R)$ 内和函数连续，且可\textbf{逐项求导}、\textbf{逐项积分}：
\keyformula{s'(x)=\sum_{n=1}^{\infty}na_nx^{n-1},\qquad \int_0^{x}s(t)\,\mathrm{d}t=\sum_{n=0}^{\infty}\frac{a_n}{n+1}x^{n+1}}
运算后\textbf{收敛半径不变}，但\textbf{端点处的收敛性可能改变}。

\corepoint{逐项求导降次、可能"丢失"常数项，逐项积分升次并添 $x$ 因子——求\textbf{和函数}正靠这两招把系数凑成可积分/可求导的形状，再回到已知展开式。逐项运算只在\textbf{开区间}内保证成立，端点要重新讨论。}

\section{函数展开成幂级数}
\subsection{泰勒级数与展开条件}

若 $f(x)$ 在 $x_0$ 的某邻域内有任意阶导数，则称
\[
\sum_{n=0}^{\infty}\frac{f^{(n)}(x_0)}{n!}(x-x_0)^n=f(x_0)+f'(x_0)(x-x_0)+\frac{f''(x_0)}{2!}(x-x_0)^2+\cdots
\]
为 $f(x)$ 在 $x_0$ 处的\textbf{泰勒级数}；$x_0=0$ 时称\textbf{麦克劳林级数}。

\textbf{定理（展开的充要条件）}\quad $f(x)$ 在该邻域内能展开成泰勒级数 $\iff$ 泰勒公式的余项满足 $\lim\limits_{n\to\infty}R_n(x)=0$。

\textbf{定理（唯一性）}\quad 若 $f(x)$ 能展开成 $x-x_0$ 的幂级数，则系数必为 $a_n=\dfrac{f^{(n)}(x_0)}{n!}$，展开式唯一。

\subsection{五个基本展开式}

\begin{center}
\begin{tabular}{@{}lll@{}}
\toprule
函数 & 展开式 & 收敛域 \\
\midrule
$\mathrm{e}^{x}$ & $\sum\limits_{n=0}^{\infty}\dfrac{x^{n}}{n!}=1+x+\dfrac{x^{2}}{2!}+\cdots$ & $(-\infty,+\infty)$ \\
$\sin x$ & $\sum\limits_{n=0}^{\infty}\dfrac{(-1)^{n}x^{2n+1}}{(2n+1)!}=x-\dfrac{x^{3}}{3!}+\dfrac{x^{5}}{5!}-\cdots$ & $(-\infty,+\infty)$ \\
$\cos x$ & $\sum\limits_{n=0}^{\infty}\dfrac{(-1)^{n}x^{2n}}{(2n)!}=1-\dfrac{x^{2}}{2!}+\dfrac{x^{4}}{4!}-\cdots$ & $(-\infty,+\infty)$ \\
$\ln(1+x)$ & $\sum\limits_{n=1}^{\infty}\dfrac{(-1)^{n-1}x^{n}}{n}=x-\dfrac{x^{2}}{2}+\dfrac{x^{3}}{3}-\cdots$ & $(-1,1]$ \\
$(1+x)^{m}$ & $1+mx+\dfrac{m(m-1)}{2!}x^{2}+\cdots$（$m$ 非非负整数） & $(-1,1)$ \\
\bottomrule
\end{tabular}
\end{center}

\keyformula{\mathrm{e}^{x}=\sum_{n=0}^{\infty}\frac{x^{n}}{n!},\qquad \frac{1}{1-x}=\sum_{n=0}^{\infty}x^{n}\ (|x|<1)}

\subsection{间接展开法}

直接按 $a_n=\frac{f^{(n)}(0)}{n!}$ 求系数（直接展开法）计算量大，通常用\textbf{间接展开法}：由已知展开式经\textbf{变量代换}、四则运算、\textbf{逐项求导}、\textbf{逐项积分}得新展开式，由唯一性保证正确。

\textbf{常用套路}：(1) 把 $f$ 化成 $\mathrm{e}^{u},\sin u,\ln(1+u),(1+u)^m,\frac{1}{1-u}$ 的形式，令 $u=\varphi(x)$；(2) 有理分式先\textbf{拆成部分分式}再用 $\frac{1}{1\mp x}$；(3) 含 $\arctan x,\arcsin x,\ln(1+x)$ 的先求导化成有理式，展开后再逐项积分。

\corepoint{间接展开最易错的是\textbf{收敛域}：代换后范围要跟着换。由 $\frac{1}{1-x}=\sum x^n\,(|x|<1)$ 得 $\frac{1}{1+x^2}=\sum(-1)^nx^{2n}$，收敛域仍是 $|x|<1$；而 $\ln(1+x)$ 逐项积分后在 $x=1$ 也收敛（得 $(-1,1]$）。\textbf{端点是否包含必须重新判}，不能照抄原级数区间。}

\section{函数的幂级数展开式的应用}

利用幂级数可作近似计算，误差由\textbf{余项}估计：(1) \textbf{求函数近似值}——把 $x$ 取小，展开到所需精度，由余项定项数；(2) \textbf{求积分近似值}——被积函数展开后\textbf{逐项积分}再截断；(3) \textbf{求极限}——用展开式代替函数，同阶项相消。

\textbf{微分方程的幂级数解法}：对 $y'=f(x,y)$，若 $f$ 在 $(x_0,y_0)$ 附近可展成幂级数，可设 $y=\sum a_n(x-x_0)^n$，代入比较同次幂系数逐次定出 $a_n$。

\corepoint{核心是\textbf{误差控制}：问"精确到 $10^{-4}$"就要由 $|r_n|\le10^{-4}$ 反解 $n$。交错级数（满足莱布尼茨条件）的余项不超过第一个被舍去项，是最快的估计。}

\section{函数项级数的一致收敛性}

\textbf{定义}\quad 设 $\sum u_n(x)$ 在 $I$ 上收敛于 $s(x)$。若对任意 $\varepsilon>0$，存在\textbf{与 $x$ 无关}的 $N$，使 $n>N$ 时对一切 $x\in I$ 有 $|r_n(x)|<\varepsilon$，则称它在 $I$ 上\textbf{一致收敛}。\textbf{魏尔斯特拉斯判别法}：若 $|u_n(x)|\le M_n$ 且 $\sum M_n$ 收敛，则 $\sum u_n(x)$ 一致收敛。一致收敛级数具有\textbf{连续性}、可\textbf{逐项积分}、可\textbf{逐项求导}（导级数须一致收敛）。

\corepoint{与普通收敛的差别只在 $N$ 是否\textbf{与 $x$ 无关}；有了它才能逐项求极限、积分、求导。幂级数在其收敛区间内部一致收敛，这正是它可逐项运算的依据。}

\section{傅里叶级数}
\subsection{三角级数与三角函数系的正交性}

形如 $\frac{a_0}{2}+\sum\limits_{n=1}^{\infty}(a_n\cos nx+b_n\sin nx)$ 的级数称\textbf{三角级数}。三角函数系 $1,\cos x,\sin x,\cos 2x,\sin 2x,\cdots$ 在 $[-\pi,\pi]$ 上\textbf{正交}：
\[
\int_{-\pi}^{\pi}\cos mx\cos nx\,\mathrm{d}x=\int_{-\pi}^{\pi}\sin mx\sin nx\,\mathrm{d}x=\int_{-\pi}^{\pi}\cos mx\sin nx\,\mathrm{d}x=0\quad(m\ne n),
\]
而 $\int_{-\pi}^{\pi}1^2\mathrm{d}x=2\pi$，$\int_{-\pi}^{\pi}\cos^2nx\,\mathrm{d}x=\int_{-\pi}^{\pi}\sin^2nx\,\mathrm{d}x=\pi$。

\subsection{函数展开成傅里叶级数}

设 $f(x)$ 是周期为 $2\pi$ 的周期函数，由正交性逐项积分得\textbf{欧拉—傅里叶公式}
\keyformula{a_n=\frac{1}{\pi}\int_{-\pi}^{\pi}f(x)\cos nx\,\mathrm{d}x\ (n=0,1,2,\cdots),\qquad b_n=\frac{1}{\pi}\int_{-\pi}^{\pi}f(x)\sin nx\,\mathrm{d}x\ (n=1,2,\cdots)}
由此系数确定的三角级数称 $f$ 的\textbf{傅里叶级数}。

\textbf{狄利克雷收敛定理}\quad 设 $f$ 是周期 $2\pi$ 的周期函数，若 (1) 一个周期内\textbf{连续或只有有限个第一类间断点}；(2) 一个周期内\textbf{至多有有限个极值点}，则其傅里叶级数收敛，且
\[
\frac{a_0}{2}+\sum_{n=1}^{\infty}(a_n\cos nx+b_n\sin nx)=
\begin{cases}
f(x), & x\ \text{为连续点},\\[3pt]
\dfrac{f(x^-)+f(x^+)}{2}, & x\ \text{为间断点}.
\end{cases}
\]

\corepoint{两个要点：(1) 系数公式的 $\frac{1}{\pi}$ 与 $a_0$ 前的 $\frac12$ 配套，写错全错；(2) 收敛定理给的是"和函数"——\textbf{间断点}处收敛于左右极限的\textbf{平均值}，端点 $x=\pm\pi$ 处收敛于 $\frac{f(-\pi^+)+f(\pi^-)}{2}$。问"展开式在某点等于什么"，往往就考这个平均值。}

\subsection{正弦级数和余弦级数}

\begin{itemize}
  \item $f$ 为\textbf{奇函数}：$a_n=0$，$b_n=\frac{2}{\pi}\int_0^{\pi}f(x)\sin nx\,\mathrm{d}x$，级数只含正弦项（\textbf{正弦级数}）。
  \item $f$ 为\textbf{偶函数}：$b_n=0$，$a_n=\frac{2}{\pi}\int_0^{\pi}f(x)\cos nx\,\mathrm{d}x$，级数只含余弦项（\textbf{余弦级数}）。
\end{itemize}

\textbf{有限区间上的展开（奇延拓与偶延拓）}：若 $f$ 只在 $[0,\pi]$ 上有定义，要展成正弦级数就作\textbf{奇延拓}（$F(-x)=-F(x)$）再按奇函数求 $b_n$；要展成余弦级数就作\textbf{偶延拓}（$F(-x)=F(x)$）再按偶函数求 $a_n$。所得级数在 $(0,\pi)$ 内收敛于 $f(x)$。

\corepoint{奇延拓的正弦级数在 $x=0$ 处收敛于 $0$（除非 $f(0)=0$）；偶延拓的余弦级数在 $x=0$ 处收敛于 $f(0^+)$。这是"端点值"类题目的标准考法；延拓只是求系数的技巧，原函数在 $[0,\pi]$ 上并未改变。}

\section{一般周期函数的傅里叶级数}
\subsection{周期为 $2l$ 的周期函数}

设 $f(x)$ 的周期为 $2l$，令 $t=\frac{\pi x}{l}$ 可化为周期 $2\pi$ 的情形，得
\keyformula{f(x)\sim\frac{a_0}{2}+\sum_{n=1}^{\infty}\left(a_n\cos\frac{n\pi x}{l}+b_n\sin\frac{n\pi x}{l}\right)}
其中
\[
a_n=\frac{1}{l}\int_{-l}^{l}f(x)\cos\frac{n\pi x}{l}\,\mathrm{d}x\ (n=0,1,2,\cdots),\qquad
b_n=\frac{1}{l}\int_{-l}^{l}f(x)\sin\frac{n\pi x}{l}\,\mathrm{d}x\ (n=1,2,\cdots).
\]
$f$ 为奇函数时 $a_n=0$，$b_n=\frac{2}{l}\int_0^{l}f\sin\frac{n\pi x}{l}\mathrm{d}x$；为偶函数时 $b_n=0$，$a_n=\frac{2}{l}\int_0^{l}f\cos\frac{n\pi x}{l}\mathrm{d}x$。狄利克雷定理同样成立。

\subsection{傅里叶级数的复数形式（简）}

由欧拉公式 $\mathrm{e}^{\mathrm{i}\theta}=\cos\theta+\mathrm{i}\sin\theta$ 得复数形式
\keyformula{f(x)\sim\sum_{n=-\infty}^{+\infty}c_n\mathrm{e}^{\mathrm{i}\frac{n\pi x}{l}},\qquad c_n=\frac{1}{2l}\int_{-l}^{l}f(x)\,\mathrm{e}^{-\mathrm{i}\frac{n\pi x}{l}}\,\mathrm{d}x\ (n=0,\pm1,\pm2,\cdots)}
周期为 $2\pi$ 时 $c_n=\frac{1}{2\pi}\int_{-\pi}^{\pi}f(x)\mathrm{e}^{-\mathrm{i}nx}\mathrm{d}x$，且 $c_0=\frac{a_0}{2}$，$c_n=\frac{a_n-\mathrm{i}b_n}{2}$，$c_{-n}=\frac{a_n+\mathrm{i}b_n}{2}$。

\corepoint{周期 $2l$ 的记法：把 $2\pi$ 情形的 $\cos nx$ 换成 $\cos\frac{n\pi x}{l}$，系数前的 $\frac{1}{\pi}$ 换成 $\frac{1}{l}$（区间长度由 $2\pi$ 变 $2l$）。\textbf{积分区间须与周期一致}（$[-l,l]$ 或任一长 $2l$ 的区间），否则系数全错。}
